\chapter{LITERATURE REVIEW}
\label{ch:literature}

% Focused review. Do NOT write a history of quantum computing. The committee is
% not made of quantum experts (see project brief), so explain just enough for a
% CS-bachelor reader to follow the methodology chapter.
% Target length: ~800-1,200 words.

\section{Classical Vector Search}
\label{sec:lit:classical}

% Cover, in order of complexity: exact brute force, FAISS flat indices,
% approximate-nearest-neighbour graphs (HNSW). Frame HNSW as the production
% standard - it is the most important comparison point on slide 11 of the
% midterm deck.

% Example sentence that exercises the citation machinery -- rewrite, but keep
% the same citation pattern so the reference list stays populated.
Exact nearest-neighbour search at scale was popularised by FAISS, a SIMD-vectorised
library for billion-scale similarity search~\cite{johnson2019faiss}. For datasets
where exact search is too expensive, the de facto approximate method is the
Hierarchical Navigable Small World graph~\cite{malkov2020hnsw}, which trades a
small recall loss for $O(\log N)$ retrieval.

[CLASSICAL VECTOR SEARCH \,--\, expand the paragraph above; add a brute-force
baseline reference if one is wanted.]

\section{Cross-Modal Embeddings}
\label{sec:lit:embeddings}

% One short subsection on CLIP. Explain contrastive learning at a high level
% and why text and image embeddings end up in the same metric space.

CLIP aligns a text encoder and an image encoder via contrastive pre-training
on 400~million image--caption pairs, producing a shared 512-dimensional space in
which cosine similarity captures cross-modal relevance~\cite{radford2021clip}.

[CLIP / CROSS-MODAL EMBEDDINGS \,--\, expand: explain ViT-B/32, mention that
we truncate the 512-dim output to 64/128/256 for benchmarking.]

\section{Quantum Similarity Primitives}
\label{sec:lit:quantum}

% Swap test (Buhrman et al.), amplitude encoding state preparation, Grover.
% Be honest about what is "speedup" and what is "amortised state preparation
% cost" - the qRAM caveat goes here, not in the methodology.

Grover's algorithm provides a quadratic speedup for unstructured
search~\cite{grover1996}. The swap test, introduced by Buhrman et
al.~\cite{buhrman2001swaptest}, estimates the squared inner product of two
quantum states with a single ancilla measurement. Both algorithms assume
that the dataset is already encoded into a quantum state; in the absence of
quantum RAM~\cite{giovannetti2008qram}, this state preparation step costs
$O(N)$ and cancels Grover's quadratic advantage at the system
level~\cite{aaronson2015readfine}.

[QUANTUM PRIMITIVES \,--\, expand: explain amplitude encoding cost, give the
formula for swap-test overlap, foreshadow the qRAM discussion in chapter~4.]

\section{Hybrid Classical--Quantum Architectures}
\label{sec:lit:hybrid}

% Brief discussion of recent hybrid retrieval pipelines that use a classical
% prefilter and a quantum reranker. Frame our hybrid_hnsw_swap_test engine as
% an instance of this pattern.

Quantum machine-learning surveys~\cite{biamonte2017qml} catalogue hybrid
classical--quantum pipelines that delegate the expensive global step to
classical hardware and reserve quantum resources for a smaller, structured
subproblem. Our \texttt{hybrid\_hnsw\_swap\_test} engine follows this pattern:
HNSW returns a small candidate pool, then the swap test reranks it on the
simulator (implemented with Qiskit~\cite{qiskit2024}).

[HYBRID ARCHITECTURES \,--\, expand: discuss other hybrid retrieval approaches
in the literature; position our work as an instance of this pattern.]
